\section{From local resolved shear to macroscopic creep}
\label{sec:bridge}

The results of Section~\ref{sec:results} bound what a strained inclusion can
do locally. To connect them with the macroscopic observation the question is
turned around: how large a stress around the inclusions would be needed to
raise the creep rate by the measured 25\%? The estimate that answers it has
three ingredients.

The first is the shape of the stress field. Outside a compact particle of
radius $a$ that is bonded to the matrix and strained, the elastic solution of
Eshelby \cite{Eshelby1957} gives a stress that falls off with distance $r$
from the particle centre as $(a/r)^{3}$. The resolved shear stress---the part
of the stress that pushes a dislocation (Section~\ref{sec:stress}), and hence
the only part that pushes a dislocation---is written in this scalar form as
$\tau(r)=\tau_m(a/r)^{3}$, with $\tau_m$ its value at the particle surface. In
a dilute dispersion with inclusion volume fraction $f$, the fraction of matrix
volume in which $|\tau|$ exceeds a given level $s$ is
$F(>s)=f\,(\tau_m/s-1)$ for $0<s<\tau_m$ (the volume of the particle itself
is subtracted), and the distribution of stress levels over the matrix volume
is given by the density $p(s)=-\mathrm{d}F/\mathrm{d}s=f\,\tau_m/s^{2}$.

The second is the response of dislocation glide to stress. At room
temperature dislocations pass obstacles by thermally activated jumps, and a
local stress $\tau$ changes the jump rate by the factor
$\exp(\pm V^{*}\tau/kT)$. Here $V^{*}$ is the activation volume of Section~\ref{sec:mobility}
\cite{CaillardMartin2003}. Because the stress around a particle
takes both signs, the two directions are weighted equally and the rate is
proportional to $\cosh(V^{*}\tau/kT)$.

The third is the averaging. The creep rate of the alloy is the rate averaged
over the matrix volume, so $\cosh(V^{*}\tau/kT)-1$ is integrated over the
density $p(s)$. The relative enhancement of the creep rate is
\begin{equation}
\frac{\langle\dot\varepsilon\rangle}{\dot\varepsilon_0}-1
 = f\,x_0\!\int_{0}^{x_0}\!\frac{\cosh u-1}{u^{2}}\,\mathrm{d}u
 = f\left[x_0\,\mathrm{Shi}(x_0)-(\cosh x_0-1)\right],
\label{eq:bridge}
\end{equation}
with $x_0=V^{*}\tau_m/kT$, where
$\mathrm{Shi}(x)=\int_0^x\sinh u\,\mathrm{d}u/u$ is the hyperbolic sine
integral. The integral is taken down to $s=0$, i.e.\ it assumes that the
stressed shells of neighbouring particles do not overlap; cutting it off where
they would (at $u_{\min}=f x_0$, where the shell of one particle reaches its
neighbours) lowers the predicted enhancement from 0.2500 to 0.2499 and raises
the required $\tau_m$ by about $0.01\%$ at every $V^{*}$ in
Table~\ref{tab:bridge}, so the tabulated values stand.
Equation~(\ref{eq:bridge}) is a model estimate of sensitivity, not an
extraction of the interface stress from the experiment: it describes the field
by a single amplitude $\tau_m$, weights positive and negative stresses
equally, and takes the spatial statistics of the $(a/r)^{3}$ field. It
matters that $\tau_m$ is the amplitude of the resolved shear at the particle
surface, not the full interface stress, of which only a part acts in any one
slip plane.

The inputs are the following. The temperature is 300~K. The inclusion content of the alloy is
0.35~wt\%, corresponding to a volume fraction $f=0.00246$, within the
0.1--0.3\% reported in Ref.~\cite{Friha2024JMMM}; Table~\ref{tab:bridge} is
evaluated at this value.
The activation volume is taken in the range $V^{*}=19$--$142\,b^{3}$, which
spans the values for aluminium alloys \cite{CaillardMartin2003}, with
$V^{*}=70\,b^{3}$ measured for Al--Mg--Si \cite{Soula2022JALCOM} as the
central case; no value of $V^{*}$ was extracted from the present simulations
(Section~\ref{sec:mobility}), so the range is a sensitivity range within the
model, not a measured input.

The output is Table~\ref{tab:bridge}. Its second column is obtained by
setting the left side of Eq.~(\ref{eq:bridge}) to 0.25 and solving for
$\tau_m$: 8.4--62.7~MPa across the range of $V^{*}$, the larger values at the
smaller activation volumes. Although the enhancement is proportional to $f$
at fixed $\tau_m$, the required $\tau_m$ is obtained by inverting a nonlinear
function and therefore does not scale as $1/f$: direct recalculation gives
9.8--73.3~MPa at $f=0.001$ and 8.1--60.3~MPa at $f=0.003$.

Three comparisons follow. If the 147~MPa of Ref.~\cite{Friha2024JMMM} were
the resolved shear stress acting on the dislocations, Eq.~(\ref{eq:bridge})
would predict enhancements from $6.8\times10^{2}$ to $2.7\times10^{46}$
(third column), at least $2.7\times10^{3}$ times the observed 0.25; 147~MPa
acting as a shear stress is incompatible with the observed enhancement. The
stress that an inclusion strained by 0.194\% actually produces is smaller.
For a bonded sphere held at that strain the exact solution \cite{Eshelby1957}
gives 41~MPa of resolved shear inside the particle and just outside it (Section~\ref{sec:sigma_r});
with this amplitude Eq.~(\ref{eq:bridge}) gives 0.045 at $V^{*}=19\,b^{3}$,
0.31 at $30\,b^{3}$ and values far above 0.25 at larger $V^{*}$ (fourth
column). The ridge of the interface cell, held at the same strain, produces
15~MPa of resolved shear just above its crest
(Section~\ref{sec:sigma_r}); with this amplitude the estimate gives
0.042 at $V^{*}=50\,b^{3}$ and 0.16 at $70\,b^{3}$, the
activation volume measured for Al--Mg--Si \cite{Soula2022JALCOM} (last
column). At a strain of 0.194\%, therefore, the stress around the inclusions
is of the order the estimate requires: the observed 0.25 is reproduced for
$V^{*}$ between about 30 and 75$\,b^{3}$, depending on which
of the two amplitudes is taken.

This is a consistency, not a confirmation, for three reasons. The estimate
is exponentially sensitive to the product $V^{*}\tau_m$: doubling $V^{*}$ at
fixed amplitude moves the prediction by two to four orders of magnitude, so
the agreement fixes little beyond the order of magnitude. The strain of
0.194\% is an input taken from the earlier estimate of the interface stress
\cite{Friha2024JMMM}, not a measured property of the inclusions; the
magnetostriction measured for bulk iron--aluminium alloys does not exceed
$10^{-4}$ \cite{Hall1959,BormioNunes2012}, twenty times less, and since the
stress is proportional to the strain an inclusion strained by $10^{-4}$ would
produce at most 2~MPa at its surface, for which Eq.~(\ref{eq:bridge}) gives
an enhancement below 0.4\% at every $V^{*}$ in the range. And the estimate
describes a field that exists while the inclusion is strained; it says
nothing about the tens of minutes for which the effect persists after the
field is removed (Section~\ref{sec:memory}). Whether the inclusions of this
alloy strain by 0.194\% in a field of 0.7~T is therefore the question on
which the elastic mechanism stands or falls, and it is an experimental one.

\begin{table}
\centering
\caption{The two-scale estimate, Eq.~(\ref{eq:bridge}), at $T=300$~K,
$f=0.00246$ and $b=2.8638$~\AA{}. Second column: the resolved
shear stress $\tau_m$ at the particle surface that reproduces the measured
$+25\%$ creep enhancement. The remaining columns give the enhancement
$\langle\dot\varepsilon\rangle/\dot\varepsilon_0-1$ that
Eq.~(\ref{eq:bridge}) predicts if $\tau_m$ were 147~MPa (the earlier
estimate), 41~MPa (the interior value of the analytical sphere held at 0.194\%, which
is also the amplitude of its inverse-cube far field; 41.45 unrounded) or 15~MPa (the maximum on the ridge axis in the interface
cell, 15.45 unrounded, Section~\ref{sec:sigma_r}). The observed value is 0.25.}
\label{tab:bridge}
\small
\begin{tabular}{lcccc}
\hline
$V^{*}$ & $\tau_m$ for $+25\%$ (MPa) & at 147~MPa & at 41~MPa & at 15~MPa \\
\hline
$19\,b^{3}$  & 62.7 & $6.8\times10^{2}$  & 0.045 & 0.004 \\
$30\,b^{3}$  & 39.7 & $3.9\times10^{6}$  & 0.31 & 0.010 \\
$50\,b^{3}$  & 23.8 & $3.9\times10^{13}$ & 17 & 0.042 \\
$70\,b^{3}$  & 17.0 & $4.8\times10^{20}$ & $1.2\times10^{3}$ & 0.16 \\
$100\,b^{3}$ & 11.9 & $2.4\times10^{31}$ & $9.3\times10^{5}$ & 1.2 \\
$142\,b^{3}$ & 8.4  & $2.7\times10^{46}$ & $1.2\times10^{10}$ & 31 \\
\hline
\end{tabular}
\end{table}

\section{Discussion}
\label{sec:discussion}

\subsection{Three stresses, and the identity of the magnetic phase}
\label{sec:phase}

Three stresses appear in this paper and must be kept apart. The first is the
147~MPa of Ref.~\cite{Friha2024JMMM}. It is an estimate of the stress at the
inclusion surface, obtained from the field through an empirical coefficient;
it is not a measured stress, not a resolved shear stress on any slip plane,
and not a result of the present calculations. If it were the shear stress on
a slip plane it would exceed the 95--105~MPa at which the lower
partner of the dislocation pair of Section~\ref{sec:thresholds} breaks away,
so the loaded cells cannot test it by comparison with a threshold; what they
test is whether a strained inclusion produces anything of that order in the
matrix. The second is what
a bonded elastic inclusion strained by 0.194\% supplies. The exact solution
for a bonded sphere \cite{Eshelby1957}---no sliding at the interface: the
atoms on both sides interact through the same potential, and no artificial
constraint is applied there---gives 41~MPa of resolved shear inside the particle,
falling as $(a/r)^{3}$ outside (Fig.~\ref{fig:rss}). The third is the field
measured in the cell for the same strain held fixed
(Section~\ref{sec:sigma_r}): 15~MPa of resolved shear just above
the crest of the ridge, falling within 80~\AA{} to $2.5\pm0.5$~MPa, the resolution of the
calculation; averaged over the width of the cell the peak is
5.0~MPa, because the ridge occupies less than half of that
width. The second and the third stress are of one order and together
bracket the amplitude entered in Eq.~(\ref{eq:bridge}). The first exceeds
them by a factor of 3.6--10 because $E\varepsilon$ is the stress of a rod held
at that strain, not of an inclusion embedded in a matrix that deforms with
it.

A separate difficulty concerns the phase itself. In the batch studied, X-ray
diffraction identifies the inclusions as Al$_{13}$Fe$_4$ (written
Fe$_4$Al$_{13}$ in Ref.~\cite{Friha2024JMMM}; the form Al$_{13}$Fe$_4$ is
used throughout here), and the magnetization of the samples shows a
hysteresis loop, i.e.\ the inclusions contain a ferromagnetic constituent
\cite{Friha2024JMMM}; earlier work in the same series gives the composition
of the inclusions as Fe$_x$Al$_{1-x}$ with $x=86$--$87$~at.\%
\cite{Skvortsov2019JAP,Nikolaev2025SciRep}. Stoichiometric Al$_{13}$Fe$_4$,
however, shows no magnetic order down to 2~K---neither in its magnetization
nor in M\"ossbauer spectra, which probe the magnetic state of the iron atoms
directly \cite{Avers2025Crystals,Albedah2015Mossbauer}: the
ideal compound does not magnetize and cannot change shape through
magnetization (which does not, by itself, guarantee a strictly zero strain of
any other origin). The measured ferromagnetism of the inclusions therefore
belongs not to the ideal Al$_{13}$Fe$_4$ lattice but to some other
constituent of them---iron-rich regions of the kind noted in
Refs.~\cite{Skvortsov2019JAP,Nikolaev2025SciRep}, for example---and it is the
strain of that constituent under field that matters for the mechanism. That
strain has not been measured for the inclusions of this alloy; the $10^{-4}$
used above is the largest value reported for bulk iron--aluminium alloys
\cite{Hall1959,BormioNunes2012}. The simulation cell represents the
inclusion by Al$_{13}$Fe$_4$ as the phase identified in this batch, with a
known structure \cite{Liu2024CIF}; the strain applied to it is taken as a
given. Measuring the strain of the actual inclusions under field is the
decisive experimental step.

\subsection{Why the mechanism cannot be simulated directly, and what MD can
measure instead}

A dislocation behaves as a line under tension. To bow it through a stressed
region of size $L$ the stress must exceed about $\mu b/L$, where $\mu$ is
the shear modulus and $b$ the Burgers vector \cite{CaillardMartin2003}. If
the whole claimed 147~MPa acted as resolved shear stress---the assumption most
favourable to the mechanism---this criterion gives $L\approx52$~nm (about
$180\,b$), i.e.\ a cell of some $2\times10^{8}$ atoms once the surrounding
matrix is included, against the roughly $10^{5}$ atoms of the cells used here.
At 41~MPa the corresponding size is about 0.2~$\mu$m, and at 15~MPa
about 0.5~$\mu$m. The mechanism therefore cannot be simulated directly at the
atomic scale. What molecular dynamics can supply is the quantities that a
description on the larger scale consumes: the stress at the interface and its
decay with distance (Section~\ref{sec:sigma_r}), the stresses at which
dislocations are set in motion or created
(Sections~\ref{sec:thresholds}--\ref{sec:mobility}), and, through
Eq.~(\ref{eq:bridge}), a check that the whole chain is self-consistent.

\subsection{The memory protocol}
\label{sec:memory}

In the experiments the field is absent during creep. The sample is held in the
field for 30~min, the field is removed, and only then is the load applied
\cite{Skvortsov2019PSS,Nikolaev2025SciRep}; what is measured is a memory of
the exposure.
Whatever the field does must therefore persist for tens of minutes after it
is gone. An elastic strain of the inclusion cannot: once the field is removed
the inclusion returns to its original shape and the stress around it
disappears within nanoseconds. The present calculations contain an inclusion held at its strain while the
matrix relaxes (the interface cell), one relaxed from it (the free cell) and
loaded cells with the inclusion free (unstrained) or held (unstrained and strained); none of them represents the
state of the material after the field has been removed. What could last
for tens of minutes is a slow rearrangement of atoms---for instance of the Mg
and Si atoms that pin the dislocations (Section~\ref{sec:mobility})---set in
motion while the field acted. Such a rearrangement has not been established
here. Supporting it would require the thermal history and the vacancy content
of the samples, an estimate of the diffusion rates, and a separate calculation
in which the field is switched on, held, switched off and the load then
applied. It is therefore stated here as one hypothesis that can be tested,
not as a result.

\subsection{Limitations}
\label{sec:limitations}

Two features of the geometry limit the transfer of these numbers to a real
particle. In the cell the inclusion is a ridge---a bar of half-elliptic cross
section running the full length of the cell along $y$
(Section~\ref{sec:cell})---rather than a compact particle. The stress outside
such a bar falls off with distance more slowly than the $(a/r)^{3}$ of a
compact particle, and its reach---some 30~\AA{} at full strength, 80~\AA{} in all
(Section~\ref{sec:sigma_r})---is set by the cross section of the ridge; it is a property of this cell, not of a micron-sized inclusion, whose
field would extend correspondingly further. Equation~(\ref{eq:bridge}) takes
its spatial statistics from the three-dimensional solution for a compact
particle and its amplitude either from that solution (41~MPa) or from the
ridge cell (15~MPa); the two differ by a factor of about
2.7, which is the effect of the geometry and is carried
through Table~\ref{tab:bridge} as the difference between its last two
columns. A two-dimensional elastic calculation for the ridge alone, with
the same strain and elastic constants but in an unbounded matrix of the same
stiffness, gives 20~MPa at the ridge surface falling to 5~MPa
within 10~\AA{}; the atomistic ridge, held rigid on its rigid layer, gives a
smaller surface value and a flatter profile. The atomistic and the continuum
description of the ridge thus agree in magnitude, and the difference from
the sphere is geometric, not an artefact of the cell.

The interatomic potential \cite{Jelinek2012} was chosen for its coverage of
Al, Fe, Mg and Si; how well it reproduces the Al/Al$_{13}$Fe$_4$ interface,
the stress needed to create dislocations, the structure of the dislocation
core and the pinning by Mg and Si has not been checked independently here, so
the absolute thresholds should be regarded as model-dependent. The residual
mismatch between the two lattices is the same in the control and strained
cells and cancels to first order in their difference; nonlinear relaxation may
prevent exact cancellation.

The thresholds of Section~\ref{sec:thresholds} are observations in a single
cell, with one set of initial atomic velocities and one loading rate; they
should be repeated with independent sets of starting velocities and at least a
second rate. The pinning result rests on two random arrangements of Mg and Si atoms,
only one of which completed the staircase, and one set of starting
velocities each; a statistical pinning threshold
would require three to five independent arrangements and velocity sets. This
is also why the pinning is reported as an observation and no activation
volume is extracted from it. The thermostat---the algorithm that holds the
cell at 300~K---acts on all mobile atoms and may itself influence the
thresholds and the temperature rise when the two dislocations meet and
annihilate; the dynamic stage should
be repeated either with the thermostat switched off after the temperature has
been established, or with it acting only on the atoms at the cell boundaries.

The field is represented only mechanically, as a strain of the inclusion.
Mechanisms in which the field acts directly on the dislocation or on the
obstacle, as in the theories of Refs.~\cite{Alshits2003,Molotskii2000}, are
outside the scope of the calculation. It is worth noting that the changes of
barrier height invoked by those theories, 0.01--0.1~eV, correspond through
$V^{*}\Delta\tau$ to 0.97--9.7~MPa at $V^{*}=70\,b^{3}$, and to
0.48--35.9~MPa over the whole range of $V^{*}$---overlapping the stresses
found here rather than coinciding with them. The stresses themselves are
computed from the atomic forces and positions (the virial expression
\cite{Thompson2022LAMMPS}) averaged over slices 4~\AA{} thick parallel to the
interface; this is an approximation whose resolution, after the six stress
components have been averaged over each slice, is $2.5\pm0.5$~MPa. Finally,
Eq.~(\ref{eq:bridge}) is a separate model relating a local stress to a change
of creep rate; it does not remove the dependence of the MD thresholds
themselves on the loading rate.

\section{Conclusions}
\label{sec:conclusions}

The stress that a strained inclusion produces in the matrix was computed in
an Al/Al$_{13}$Fe$_4$ interface cell of 91\,428 atoms for an elongation of
the inclusion by 0.194\% along the field direction at constant volume, the
strain corresponding to the earlier estimate of 147~MPa, with the inclusion
held at that strain. The resolved shear stress in the matrix reaches
15~MPa just above the inclusion and falls to the far-field level
of $2.5\pm0.5$~MPa within 80~\AA{}; averaged over the width of the cell the
peak is 5.0~MPa. The analytical solution for a compact particle
held at the same strain gives 41~MPa inside it, falling as the inverse
cube of the distance, so that for a micron-sized inclusion the same field
extends over a fraction of a micron. An inclusion that is not held relaxes
most of the strain.

In the loaded cells, a pre-existing pair of dislocations next to the ridge
is torn apart at an applied shear of 95--105~MPa with the inclusion free to relax; with the inclusion held, the strained and the unstrained cell coincide within one frame of the ramp, 9~MPa, at every stage; no new dislocation forms at the interface up to the end of the ramp at 400~MPa in the unstrained cell; and a random
arrangement of Mg and Si atoms holds a dislocation through 75~MPa. The first picoseconds of every ramp, at zero applied stress, answer the question whether the field of the strained ridge alone moves a dislocation placed next to it: the same partner moves in the same way with the ridge strained and unstrained, driven by the coherency field of the interface, not by the strain of the ridge.
These values are specific to the cells, the loading rate and the potential
and are not material constants.

The two-scale estimate, Eq.~(\ref{eq:bridge}), requires a resolved shear
stress of 8.4--62.7~MPa at the inclusion surface to raise the creep rate by
25\% at $f=0.00246$ and $V^{*}=19$--$142\,b^{3}$. The computed amplitudes,
41~MPa for the particle and 15~MPa for the ridge, reproduce the
observed enhancement for $V^{*}$ of 30--75$\,b^{3}$, a range
that contains the value measured for Al--Mg--Si; 147~MPa acting as resolved
shear would exceed the observed enhancement by a factor of at least
$2.7\times10^{3}$. The elastic mechanism is therefore quantitatively viable
if, and only if, the inclusions strain by about 0.2\% in the field: the
magnetostriction measured for bulk iron--aluminium alloys is twenty times
smaller and would give an enhancement below 0.4\%.

Because the field is off during the creep test, whatever it does must persist
for tens of minutes; an elastic strain cannot, and a slow rearrangement of
atoms is indicated as the remaining possibility---a hypothesis not tested by
the present calculations. The decisive next steps are the identification of
the ferromagnetic constituent of the inclusions, a direct measurement of their
strain under field, and a calculation in which the field is switched on, held,
switched off and the load then applied.
